\section*{Capítulo \ref{ch:b2:catalytic-cycles} --- Catálisis organometálica: etapas elementales y ciclos}

\begin{solution}{exo:b2:catalytic-cycles:1}
Antes: Ir(I), $d^8$, $8 + 4 \times 2 = 16$ electrones. Después: Ir(III), $d^6$, $6 + 6 \times 2 = 18$
electrones, octaédrico.
\end{solution}

\begin{solution}{exo:b2:catalytic-cycles:2}
La primera es una \omterm{def:b2:catalytic-cycles:oxidative-addition}{eliminación reductora} (de Pt(IV) a Pt(II), con formación de etano a partir de dos metilos).
La segunda es un \omterm{def:b2:catalytic-cycles:ligand-exchange}{intercambio de ligandos} (una fosfina sustituida por eteno) seguido de una
\omterm{def:b2:catalytic-cycles:insertion}{inserción migratoria} del eteno en el enlace Pd--Ph.
\end{solution}

\begin{solution}{exo:b2:catalytic-cycles:3}
TON $= 15.0/0.020 = 750$; TOF $= 750/3.0 = \qty{250}{h^{-1}}$.
\end{solution}

\begin{solution}{exo:b2:catalytic-cycles:4}
\begin{align*}
  &\mathrm{PdL_2 + ArBr \to ArPdBrL_2}\\
  &\mathrm{ArPdBrL_2 + Ar'B(OH)_3^- \to ArPdAr'L_2 + Br^- + B(OH)_3}\\
  &\mathrm{ArPdAr'L_2 \to Ar{-}Ar' + PdL_2}
\end{align*}
Suma: $\mathrm{ArBr + Ar'B(OH)_3^- \to Ar{-}Ar' + Br^- + B(OH)_3}$; todas las especies de paladio
se cancelan.
\end{solution}

\begin{solution}{exo:b2:catalytic-cycles:5}
El etilo y el isopropilo (tienen hidrógenos en $\beta$). El metilo no tiene carbono $\beta$; el carbono $\beta$
del neopentilo no lleva hidrógenos; el carbono $\beta$ del bencilo es un carbono \omterm{def:b2:huckel:aromatic}{aromático} del anillo sin
hidrógeno.
\end{solution}

\begin{solution}{exo:b2:catalytic-cycles:6}
(\textit{E})-3-Fenilpropenoato de metilo, \ce{Ph-CH=CH-COOCH3}, con el fenilo unido al
carbono terminal: producto trans.
\end{solution}

\begin{solution}{exo:b2:catalytic-cycles:7}
Butanal, \ce{CH3CH2CH2CHO} (lineal: el metal se une al carbono terminal, y el hidruro,
al interno), y 2-metilpropanal, \ce{(CH3)2CHCHO} (ramificado: el metal sobre el carbono
interno).
\end{solution}

\begin{solution}{exo:b2:catalytic-cycles:8}
\ce{[RhCl(PPh3)3]} tiene 16 electrones: adicionar \ce{H2} (+2) daría 18 con seis
posiciones de coordinación, demasiado congestionado con tres fosfinas voluminosas. \ce{[RhCl(PPh3)2]} tiene 14
y una posición libre: la \omterm{def:b2:catalytic-cycles:oxidative-addition}{adición oxidante} da con facilidad un complejo de 16 electrones.
\end{solution}

\begin{solution}{exo:b2:catalytic-cycles:9}
El ácido borónico es un mal nucleófilo; la base adiciona hidróxido (o un alcóxido) al boro,
dando un borato \ce{ArB(OH)3-}, cuyo grupo arilo es más rico en electrones y se transfiere al
paladio.
\end{solution}

\begin{solution}{exo:b2:catalytic-cycles:10}
Un complejo de 20 electrones no es razonable. El complejo debe perder primero un ligando (14 e), luego
adicionar \ce{H2} (16 e), unir el alqueno (18 e), insertarlo en un enlace M--H (16 e) y eliminar el
alcano (14 e), antes de recuperar el ligando o empezar de nuevo.
\end{solution}

\begin{solution}{exo:b2:catalytic-cycles:11}
Velocidad inicial $n_{\max}/\tau = \qty{0.25}{mmol/min}$; TOF inicial $= 0.25/0.010 =
\qty{25}{min^{-1}} = \qty{1.5e3}{h^{-1}}$. TON final $= 10.0/0.010 = 1.0 \times 10^3$.
\end{solution}

\begin{solution}{exo:b2:catalytic-cycles:12}
4-Bromoanisol con ácido fenilborónico (o bromobenceno con ácido 4-metoxifenilborónico),
un catalizador de paladio(0) con fosfina y una base. Ciclo: \omterm{def:b2:catalytic-cycles:oxidative-addition}{adición oxidante} del
bromuro de arilo, \omterm{def:b2:catalytic-cycles:transmetalation}{transmetalación} desde el borato y \omterm{def:b2:catalytic-cycles:oxidative-addition}{eliminación reductora} del
4-metoxibifenilo.
\end{solution}

\begin{solution}{pb:b2:catalytic-cycles:1}
\textbf{1.} Pd(II), $d^8$.
\textbf{2.} Pd(0), $d^{10}$, $10 + 2 \times 2 = 14$ electrones.
\textbf{3.} Con 14 electrones puede aceptar dos pares más: tiene posiciones vacantes.
\textbf{4.} Los complejos fosfínicos de paladio(0) se oxidan con el aire, y las fosfinas se
oxidan a óxidos de fosfina.
\textbf{5.} Reduce el paladio(II) a paladio(0) (oxidándose ella misma) y se une al
paladio(0), manteniéndolo en disolución.
\textbf{6.} El \omterm{def:b2:catalytic-cycles:cycle}{precatalizador}.
\textbf{7.} \ce{Pd(PPh3)2 + CH3C6H4Br -> [Pd(C6H4CH3)Br(PPh3)2]}: Pd(II), $d^8$, $8 + 4 \times 2 =
16$ electrones.
\textbf{8.} Lo convierte en el borato \ce{PhB(OH)3-}, que transfiere su grupo fenilo.
\textbf{9.} \ce{[Pd(Ar)Br(PPh3)2] + PhB(OH)3- -> [Pd(Ar)(Ph)(PPh3)2] + Br- + B(OH)3}.
\textbf{10.} La \omterm{def:b2:catalytic-cycles:oxidative-addition}{eliminación reductora} une dos ligandos en cis; los arilos en trans deben isomerizarse primero.
\textbf{11.} \ce{[Pd(Ar)(Ph)(PPh3)2] -> Ar-Ph + Pd(PPh3)2}: paladio(0), 14 electrones, listo
para otra vuelta.
\textbf{12.} \ce{CH3C6H4Br + PhB(OH)3- -> CH3C6H4-Ph + Br- + B(OH)3}.
\textbf{13.} El 4-bromotolueno (\qty{10.0}{mmol}, frente a \qty{12.0}{mmol} de ácido borónico).
\textbf{14.} $1.51/168.24 = \qty{8.98e-3}{mol}$, \qty{8.98}{mmol}.
\textbf{15.} $8.98/10.0 = 89.8~\%$.
\textbf{16.} $0.050/10.0 = 0.50$~mol~\%.
\textbf{17.} $8.98/0.050 = 180$.
\textbf{18.} $180/4.0 = \qty{45}{h^{-1}}$.
\textbf{19.} Dos grupos fenilo transferidos al mismo paladio (homoacoplamiento del
ácido borónico, favorecido por trazas de oxígeno), y eliminados después juntos.
\textbf{20.} Los grupos arilo no tienen ningún hidrógeno en $\beta$ capaz de alcanzar el metal.
\textbf{21.} La \omterm{def:b2:catalytic-cycles:oxidative-addition}{adición oxidante}: el enlace \ce{C-Cl} es más fuerte que el \ce{C-Br}.
\textbf{22.} Termina como partículas de paladio(0) o como sales en los residuos; es caro y
tóxico, y se recupera.
\textbf{23.} Parte del ácido borónico se pierde en reacciones secundarias (homoacoplamiento, pérdida del boro
del anillo); el exceso garantiza que el bromuro siga siendo el limitante.
\textbf{24.} El paladio completó $\boldsymbol{\approx \num{1.8e2}}$ ciclos.
\end{solution}
